\chapter{L’Esprit saint et la Trinité}\label{ch:esprit_trinite}

\section{Une seule mention possible de l’Esprit saint}

Le \nnum{249} est le seul texte de l’ensemble C qui puisse mentionner l’Esprit saint. Il s’agit d’un petit fragment de plaque (20 × 23,5 cm selon l’EDB), rattaché par l’EDB au cimetière de Nicomède, sur la via Nomentana. Les deux bases en donnent des transcriptions proches :

\begin{quote}\small
EDCS : \lat{[- - -]tus est qui ac[cepti sanc]tum spiritum [XVI Kal(endas)] Mai(as) Merobaud[e II con(sule)]}\\
EDB : \lat{tus est qui ac[cepit sanc]tum spirituṃ [XVI ? kal(endas)] mai(as) Merobaud[e consule]}
\end{quote}

Le mot \emph{spiritum} est lu, la dernière lettre étant pointée par l’EDB. De l’adjectif ne subsiste que la finale \emph{-tum} ; la restitution \emph{[sanc]tum} est vraisemblable, mais d’autres adjectifs sont théoriquement possibles. Le verbe est restitué : \emph{ac[cepti]} dans l’EDCS, forme qui n’est pas grammaticale en l’état et que l’on ne corrige pas ici, \emph{ac[cepit]} dans l’EDB. Le quantième est restitué (l’EDB y ajoute un point d’interrogation). L’année repose sur la lecture \emph{Merobaud[- - -]} : l’itération \emph{II} est une restitution de l’EDCS, que l’EDB ne reprend pas. Les deux bases rapportent néanmoins le texte à 383, année du deuxième consulat de Merobaude, où celui-ci est nommé en premier ; la place du nom dans la formule s’accorde avec cette date.

Si la restitution est retenue, le texte dit d’un défunt « qu’il reçut l’Esprit saint », peu avant sa mort ou à une date indiquée. L’expression peut désigner l’initiation chrétienne, et plus précisément le rite postbaptismal que les textes contemporains associent au don de l’Esprit. Mais rien dans le fragment ne permet de préciser ce rite, ni de dater l’événement par rapport au décès. Une coïncidence de calendrier mérite d’être signalée, pour en marquer les limites. Selon le comput alexandrin, Pâques tombait en 383 le 9 avril, et le 16 avril (\emph{a. d. XVI Kal. Mai.}) était le dimanche suivant (calcul vérifié par le jour julien). Mais le quantième est entièrement restitué, l’EDB l’assortit d’un point d’interrogation, et l’Église de Rome n’appliquait pas nécessairement le comput alexandrin à cette date : aucune conclusion liturgique ne peut être tirée de ce rapprochement. Le texte est retenu dans le profil large et exclu du profil restrictif (chapitre~\ref{ch:typologie}).

Ailleurs, \emph{spiritus} désigne l’esprit du défunt : \emph{spiritus a carne recedens}, « l’esprit se retirant de la chair » (\nnum{73}), \emph{cuius spiritus in luce domini susceptus est} (\nnum{400}). Ces emplois relèvent de l’anthropologie, la séparation de l’esprit et du corps, non de la pneumatologie.

\section{\emph{Confirmans trinitatem}}

L’épitaphe de Quintilianus (\nnum{474}, 403) est conservée sur sa moitié gauche ; le texte est identique dans les deux bases :

\begin{quote}\small
\lat{Quintilianus homo dei [- - -] / confirmans trinitatem [- - -] / amans castitatem [- - -] / respuens mundum [- - -] / requiescet una cum [- - -] / qui vixit ann(os) p(lus) m(inus) X[- - -] / Theodosio Aug(usto) et R[umorido v(iro) c(larissimo) conss(ulibus)]}
\end{quote}

Le mot \emph{trinitatem} est lu. Les trois participes forment un éloge ascétique : affirmation de la foi, amour de la chasteté, mépris du monde. Le parallélisme invite à comprendre \emph{confirmans trinitatem} comme une qualité religieuse du défunt, son adhésion déclarée à la foi trinitaire, plutôt que comme une formule dogmatique autonome. Le participe \emph{confirmans} peut signifier « affirmant », « professant » ou « fortifiant » ; la mutilation de chaque ligne à droite interdit de savoir si un complément suivait. Le texte ne nomme ni le Père, ni le Fils, ni l’Esprit ; le profil « trinitaire explicite » du catalogue de travail est remplacé par « mention de la Trinité ».

L’année 403 repose sur \emph{Theodosio Aug(usto)}, lu, et sur l’initiale \emph{R}, complétée en \emph{R[umorido]}. Entre 340 et 403, seul Rumoridus est le collègue d’un Théodose Auguste avec un nom commençant par R ; ce raisonnement reste à confirmer sur la CLRE. Selon l’EDB, la pierre, aujourd’hui perdue, provient du cimetière à ciel ouvert proche de la basilique Saint-Pierre.

\section{Contexte et limites de l’interprétation}

Le contexte de ces deux textes peut être rappelé sans qu’il soit possible d’en tirer une explication. Le \nnum{474} est postérieur d’une vingtaine d’années au concile de Constantinople (381) et à la législation de Théodose en faveur de la foi de Nicée, dont la documentation n’a pas été reprise dans ce volume. Que l’unique mention de la Trinité du corpus soit l’une de ses dernières inscriptions est un fait. Qu’elle traduise la diffusion d’un langage doctrinal dans l’épigraphie privée est une hypothèse que ce seul texte ne permet pas de vérifier : il faudrait la confronter aux épitaphes de la première moitié du \textsc{v}\textsuperscript{e} siècle, qui sortent du cadre de ce volume.

Dans l’état du corpus, l’Esprit saint est nommé au plus une fois, sur restitution partielle ; la Trinité une fois, sans énumération des personnes.
